\ExerciseSection

\begin{enumerate}
\def\labelenumi{\arabic{enumi})}
\item
  Let X be a topological space and Y a metric space. Show that the set of all mappings of X into Y whose oscillation (Chapter IX, § 2, no. 3) at every point of X is \(\leqslant\alpha\) (where \(\alpha\) is a given real number \textgreater0) is closed in the space \(\mathcal{F}_{u}(\mathbf{X};\mathbf{Y})\) .
\item
  Let X be a set; show that the mapping \(u \to \sup_{x \in X} u(x)\) of \(\mathcal{B}(X; \mathbf{R})\) into R is continuous.
\item
  Let \(\mathbf{X}\) be a metric space, \(d\) the metric on \(\mathbf{X}\) . For each \(x \in \mathbf{X}\) let \(d_x\) be the real-valued function \(y \to d(x, y)\) which is continuous on \(\mathbf{X}\) . Show that the mapping \(x \to d_x\) is an isometry of \(\mathbf{X}\) onto a subspace of \(\mathcal{C}_u(\mathbf{X}; \mathbf{R})\) {[}endowed with the pseudometric \(\delta(u, v) = \sup_{x \in \mathbf{X}} |u(x) - v(x)|\) {]}.
\item
  A completely regular space X is such that the metrizable space \(\mathcal{C}_{u}(\mathbf{X};\mathbf{R})\) is of countable type if (and only if) X is compact and metrizable. {[}By considering the Stone-Čech compactification of X (Chapter IX, § 1, Exercise 7) show that X must be metrizable, and then observe that the space \(\mathcal{C}_{u}(\mathbf{Z};\mathbf{R})\) is not of countable type.{]}
\item
  Let X be a completely regular space.
\end{enumerate}

\begin{enumerate}
\def\labelenumi{\alph{enumi})}
\tightlist
\item
  If every point of the space \(\mathcal{C}_{c}(\mathbf{X};\mathbf{R})\) has a countable fundamental system of neighbourhoods, then there exists an increasing sequence \((\mathbf{K}_{n})\) of compact subsets of X such that every compact subset of X is contained in some \(K_{n}\) ; and \(\mathcal{C}_{c}(\mathbf{X};\mathbf{Y})\) is then metrizable for any metrizable space Y.
\item
  \(\mathcal{C}_{c}(\mathbf{X};\mathbf{R})\) is metrizable and of countable type if and only if all the compact subspaces \(K_{n}\) are metrizable (use Exercise 4). For every metrizable space Y of countable type, \(\mathcal{C}_{c}(\mathbf{X};\mathbf{Y})\) is then metrizable and of countable type.
\end{enumerate}

\begin{enumerate}
\def\labelenumi{\arabic{enumi})}
\setcounter{enumi}{5}
\tightlist
\item
  \begin{enumerate}
  \def\labelenumii{\alph{enumii})}
  \tightlist
  \item
    Let X be a topological space, Y a metric space, d the metric on Y, n an integer \textgreater0 and S a subset of the space \(X^{n} \times Y^{n} \times R_{+}^{*}\) . For each point \(z = ((x_{i}), (y_{i}), r) \in \mathbf{S}\) , let \(U_{z}\) be the open subset of \(\mathcal{C}_{s}(X; Y)\) consisting of continuous mappings f of X into Y such that \(d(f(x_i), y_i) < r\) for \(1 \leqslant i \leqslant n\) . Show that, if D is a dense subset of S, we have \(\bigcup_{z \in \mathbf{S}} \mathrm{U}_z = \bigcup_{z \in \mathbf{D}} \mathrm{U}_z\) .
  \end{enumerate}
\end{enumerate}

\begin{enumerate}
\def\labelenumi{\alph{enumi})}
\setcounter{enumi}{1}
\tightlist
\item
  Deduce from a) that if the topologies of X and Y have countable bases, then every subspace of \(\mathcal{C}_{s}(\mathbf{X};\mathbf{Y})\) is a Lindelöf space (Chapter I, § 9, Exercise 15) and is therefore paracompact (Chapter IX, § 4, Exercise 23).
\end{enumerate}

\begin{enumerate}
\def\labelenumi{\arabic{enumi})}
\setcounter{enumi}{6}
\item
  Let X be a space which has a countable dense subset, let Y be a Hausdorff space in which every point has a countable fundamental system of neighbourhoods (resp. a metrizable space) and let H be a subset of C(X; Y). Show that if C is a topology on H which is finer than the topology of pointwise convergence in X and with respect to which H is compact, then every point of H has a countable fundamental system of neighbourhoods in the topology C (resp. C is metrizable). (Note that if D is any countable dense subset of X, then the topology C is finer than the topology of pointwise convergence in D and that the latter topology is Hausdorff.)
\item
  \begin{enumerate}
  \def\labelenumii{\alph{enumii})}
  \tightlist
  \item
    Let \(f\) be the continuous mapping \((x, y) \to xy\) of \(\mathbf{R} \times \mathbf{R}\) into \(\mathbf{R}\) . Show that the mapping \(x \to f(x, .)\) of \(\mathbf{R}\) into \(\mathcal{C}_u(\mathbf{R}; \mathbf{R})\) is not continuous.
  \end{enumerate}
\end{enumerate}

\begin{enumerate}
\def\labelenumi{\alph{enumi})}
\setcounter{enumi}{1}
\tightlist
\item
  Let X, Y be two topological spaces, let Z be a uniform space and let \(f: \mathbf{X} \times \mathbf{Y} \to \mathbf{Z}\) be a mapping. Show that if \(f(x, .)\) is continuous on Y for all \(x \in \mathbf{X}\) and if the mapping \(x \to f(x, .)\) of X into \(\mathcal{C}_u(Y; \mathbf{Z})\) is continuous, then \(f\) is continuous on \(\mathbf{X} \times \mathbf{Y}\) .
\end{enumerate}

\begin{enumerate}
\def\labelenumi{\arabic{enumi})}
\setcounter{enumi}{8}
\tightlist
\item
  \begin{enumerate}
  \def\labelenumii{\alph{enumii})}
  \tightlist
  \item
    Let X, Y, Z be three topological spaces. Suppose that X and Y are Hausdorff and that every point of each of these spaces has a countable fundamental system of neighbourhoods. Let f be a mapping of \(X \times Y\) into Z such that \(f(x, .)\) is continuous on Y for each \(x \in X\) and such that the mapping \(x \to f(x, .)\) of X into \(\mathcal{C}_{c}(Y; Z)\) is continuous. Show that f is continuous on \(X \times Y\) (cf.~§1, no. 6, Theorem 2, Corollary 3).
  \end{enumerate}
\end{enumerate}

\begin{enumerate}
\def\labelenumi{\alph{enumi})}
\setcounter{enumi}{1}
\tightlist
\item
  Show that the conclusion of a) holds good when the hypothesis on X is replaced by the hypothesis that X is locally compact and Z is a uniform space (use § 2, Exercise 2, and Ascoli's theorem) {[}cf.~Exercise II b){]}.
\end{enumerate}

¶ 10) Let X, Y be two topological spaces. For each subset A of X and each subset B of Y let \(T(A, B)\) denote the set of all \(u \in \mathcal{C}(X; Y)\) such that \(u(A) \subset B\) .

Let \(\mathfrak{U} = (\mathbf{U}_{\alpha})\) be an open covering of \(\mathbf{X}\) . Let \(\mathfrak{G}_{\mathfrak{U}}\) denote the topology on \(\mathcal{C}(\mathbf{X};\mathbf{Y})\) generated by the sets \(T(\mathbf{F},\mathbf{V})\) , where \(\mathbf{V}\) runs through the set of all open sets of \(\mathbf{Y}\) and \(\mathbf{F}\) runs through the set of closed subsets of \(\mathbf{X}\) contained in at least one \(\mathbf{U}_{\alpha}\) .

\begin{enumerate}
\def\labelenumi{\alph{enumi})}
\item
  Show that \(\mathfrak{G}_{\mathfrak{U}}\) is finer than the compact-open topology (use Theorem 3 of Chapter IX, § 4, no. 3). If \(\mathbf{X}\) is regular, the mapping \((u, x) \to u(x)\) of \(\mathcal{C}(\mathbf{X}; \mathbf{Y}) \times \mathbf{X}\) into \(\mathbf{Y}\) is continuous when \(\mathcal{C}(\mathbf{X}, \mathbf{Y})\) carries the topology \(\mathfrak{G}_{\mathfrak{U}}\) .
\item
  Let \(\mathfrak{C}\) be a topology on \(\mathcal{C}(\mathbf{X};\mathbf{Y})\) with respect to which the mapping \((u,x)\to u(x)\) of \(\mathcal{C}(\mathbf{X};\mathbf{Y})\times \mathbf{X}\) into \(\mathbf{Y}\) is continuous. Let \(u_0\) be a continuous mapping of \(\mathbf{X}\) into \(\mathbf{Y}\) , let \(x_0\) be a point of \(\mathbf{X}\) and let \(\mathbf{V}\) be an open set in \(\mathbf{Y}\) which contains \(u_0(x_0)\) . Show that there exists an open neighbourhood \(\mathbf{U}\) of \(x_0\) in \(\mathbf{X}\) such that \(\pmb{T}(\mathrm{U},\mathrm{V})\) is a neighbourhood of \(u_0\) with respect to \(\mathfrak{C}\) .
\item
  Suppose that X is completely regular. Show that if, among the topologies \(\mathcal{C}\) on \(\mathcal{C}(\mathbf{X};\mathbf{R})\) for which \((u,x)\to u(x)\) is continuous, there is a topology \(\mathcal{C}_0\) coarser than all the others, then \(\mathcal{C}_0\) must be the compact-open topology. {[}Let \(T(\mathrm{U},\mathrm{V})\) be a neighbourhood of 0 in \(\mathcal{C}(\mathbf{X};\mathbf{R})\) with respect to \(\mathcal{C}_0\) . Let \((\mathrm{W}_{\alpha})\) be any open covering of \(\overline{\mathbf{U}}\) ; let \(\mathfrak{B}\) be the covering of X formed by the sets \(\mathrm{W}_{\alpha}\) and \(\mathbb{C}\overline{\mathbf{U}}\) . Using a) and arguing by contradiction, show that a finite number of the \(\mathrm{W}_{\alpha}\) cover \(\overline{\mathbf{U}}\) .{]}
\item
  Show that, if X is completely regular but not locally compact, the mapping \((u, x) \to u(x)\) of \(\mathcal{C}_{c}(\mathbf{X}; \mathbf{R}) \times \mathbf{X}\) into R is not continuous at every point. {[}Consider a point \(x_{0}\) which has no compact neighbourhood and argue by contradiction, using b).{]}
\end{enumerate}

¶ 11) Let X, Y be two non-empty topological spaces whose product \(X \times Y\) is normal. Let C be a topology on the set \(\mathcal{C}(X; \mathbf{R})\) such that, for every continuous real-valued function f on \(X \times Y\) , the mapping \(y \to f(., y)\) of Y into \(\mathcal{C}(X; \mathbf{R})\) is continuous.

\begin{enumerate}
\def\labelenumi{\alph{enumi})}
\item
  Let \(y_0\) be a limit of a sequence \((z_n)\) of points of Y all distinct from \(y_0\) , let A be a countably infinite closed subset of X, all of whose points are isolated, and let I be a bounded open interval in R. Show that, with respect to the topology \(\mathfrak{G}\) , the set \(T(\mathbf{A}, \mathbf{I})\) (Exercise 10) has no interior point. {[}If \(u_0 \in T(\mathbf{A}, \mathbf{I})\) , construct a continuous mapping \(f: \mathbf{X} \times \mathbf{Y} \to \mathbf{R}\) such that \(f(., y_0) = u_0\) and that \(f(., z_n) \notin T(\mathbf{A}, \mathbf{I})\) for all \(z_n \neq y_0\) .{]}
\item
  Hence show that if, in addition, the mapping \((u, x) \to u(x)\) of \(\mathcal{C}(\mathbf{X}; \mathbf{R}) \times \mathbf{X}\) into R is continuous with respect to the topology C, if X is locally paracompact (Chapter IX, § 4, Exercise 27), and if there exists in X a sequence of points \((x_{n})\) which converges to a point distinct from all the \(x_{n}\) , then, X must be locally compact {[}use Exercise 10 b) and Chapter IX, § 4, Exercise 25 c){]}.
\item
  Deduce from b) that if X is metrizable but not locally compact, then there are points of \(\mathcal{C}_{c}(\mathbf{X};\mathbf{R})\) which have no countable fundamental system of neighbourhoods {[}use Exercise 9 a){]}. Give a direct proof, using Exercise 5 a).
\end{enumerate}

\begin{enumerate}
\def\labelenumi{\arabic{enumi})}
\setcounter{enumi}{11}
\tightlist
\item
  Let X be a topological space, Y and Z two uniform spaces, S a set of subsets of X and T a set of subsets of Y.
\end{enumerate}

\begin{enumerate}
\def\labelenumi{\alph{enumi})}
\item
  Let \(u_0 \in \mathcal{C}(\mathbf{X}; \mathbf{Y})\) . Show that if, for each \(\mathbf{A} \in \mathfrak{S}\) , \(u_0(\mathbf{A})\) is contained in some set of \(\mathfrak{S}\) , then the mapping \(v \to v \circ u_0\) of \(\mathcal{C}_{\mathfrak{T}}(\mathbf{Y}; \mathbf{Z})\) into \(\mathcal{C}_{\mathfrak{S}}(\mathbf{X}; \mathbf{Z})\) is continuous.
\item
  Let H be a subspace of \(\mathcal{C}_{\mathfrak{S}}(\mathrm{X};\mathrm{Y})\) and let \(v_{0}\in\mathcal{C}(\mathrm{Y};\mathrm{Z})\) . Show that if \(v_{0}\) is uniformly continuous on \(\mathrm{H}(\mathrm{A})\) for all \(A\in S\) , then the mapping \(u\to v_{0}\circ u\) of H into \(\mathcal{C}_{\mathfrak{S}}(\mathrm{X};\mathrm{Z})\) is continuous.
\item
  Let \(u_0 \in \mathcal{C}(\mathbf{X}; \mathbf{Y})\) and let \(v_0 \in \mathcal{C}(\mathbf{Y}; \mathbf{Z})\) . Suppose that, for each \(\mathbf{A} \in \mathfrak{S}\) , there exists \(\mathbf{B} \in \mathfrak{T}\) and an entourage \(\mathbf{V}\) of \(\mathbf{Y}\) satisfying the conditions: (i) \(\mathbf{V}(u_0(\mathbf{A})) \subset \mathbf{B}\) , (ii) \(v_0\) is uniformly continuous on \(\mathbf{B}\) .
\end{enumerate}

Then the mapping \((u, v) \to v \circ u\) of \(\mathcal{C}_{\mathfrak{S}}(\mathbf{X}; \mathbf{Y}) \times \mathcal{C}_{\mathfrak{T}}(\mathbf{Y}; \mathbf{Z})\) into \(\mathcal{C}_{\mathfrak{S}}(\mathbf{X}; \mathbf{Z})\) is continuous at \((u_0, v_0)\) . In particular, the mapping \((u, v) \to v \circ u\) of \(\mathcal{C}_u(\mathbf{X}; \mathbf{Y}) \times \mathcal{C}_u(\mathbf{Y}; \mathbf{Z})\) into \(\mathcal{C}_u(\mathbf{X}; \mathbf{Z})\) is continuous at every point \((u_0, v_0)\) such that \(v_0\) is uniformly continuous on \(\mathbf{Y}\) .

\begin{enumerate}
\def\labelenumi{\alph{enumi})}
\setcounter{enumi}{3}
\tightlist
\item
  Let \(v_0\) be the homeomorphism \(x \to x^3\) of \(\mathbf{R}\) onto itself. Show that the mapping \(u \to v_0 \circ u\) of \(\mathcal{C}_u(\mathbf{R}; \mathbf{R})\) into itself is not continuous at every point.
\end{enumerate}

¶ 13) Let X, Y, Z be three Hausdorff topological spaces.

\begin{enumerate}
\def\labelenumi{\alph{enumi})}
\item
  Show that for all \(u_0 \in \mathcal{C}(\mathbf{X}; \mathbf{Y})\) and all \(v_0 \in \mathcal{C}(\mathbf{Y}; \mathbf{Z})\) , the mappings \(v \to v \circ u_0\) of \(\mathcal{C}_c(\mathbf{Y}; \mathbf{Z})\) into \(\mathcal{C}_c(\mathbf{X}; \mathbf{Z})\) and \(u \to v_0 \circ u\) of \(\mathcal{C}_c(\mathbf{X}; \mathbf{Y})\) into \(\mathcal{C}_c(\mathbf{X}; \mathbf{Z})\) are continuous.
\item
  Suppose that every point of Y (resp. Z) has a countable fundamental system of neighbourhoods and that X has a countable dense subset. Let H be a compact subset of \(\mathcal{C}_{c}(\mathbf{X};\mathbf{Y})\) . Show that the mapping \((u,v)\to v\circ u\) of \(\mathbf{H}\times\mathcal{C}_{c}(\mathbf{Y};\mathbf{Z})\) into \(\mathcal{C}_{c}(\mathbf{X};\mathbf{Z})\) is continuous. {[}Start by showing that if K is any compact subset of X, then \(\mathbf{H}(\mathbf{K})\) is a compact subset of Y, by using Exercises 7 and 9 a).{]}
\item
  Show that the mapping \((u, v) \to v \circ u\) of \(\mathcal{C}_c(\mathbf{Q}; \mathbf{Q}) \times \mathcal{C}_c(\mathbf{Q}; \mathbf{Q})\) into \(\mathcal{C}_c(\mathbf{Q}; \mathbf{Q})\) is not continuous at any point.
\end{enumerate}

\begin{enumerate}
\def\labelenumi{\arabic{enumi})}
\setcounter{enumi}{13}
\tightlist
\item
  Let \(\Gamma\) be the group of all homeomorphisms of the real plane \(\mathbf{R}^2\) , endowed with the topology of pointwise convergence. Show that the mapping \((u, v) \to v \circ u\) of \(\Gamma \times \Gamma\) into \(\Gamma\) is not continuous. {[}Consider a sequence of homeomorphisms \((v_n)\) such that \(v_n\) leaves fixed every point \((x, y)\) for which \(y \notin [1/(n + 1), 1/n]\) , and such that the restriction of \(v_{n}\) to the line \(y = 2/(2n + 1)\) is the translation \((x, y) \to (x + 1, y)\) ; on the other hand, consider the sequence \((u_{n})\) of translations
\end{enumerate}

\[
(x, y) \rightarrow (x, y + 2 / (2 n + 1)).
\]
{]}

\begin{enumerate}
\def\labelenumi{\arabic{enumi})}
\setcounter{enumi}{14}
\tightlist
\item
  \begin{enumerate}
  \def\labelenumii{\alph{enumii})}
  \tightlist
  \item
    Let X be a uniform space, let S be a set of subsets of X, let \(\Gamma\) be the group of all homeomorphisms of X, and let \(u_{0} \in \Gamma\) . Suppose that for each \(A \in S\) there exist \(B \in S\) and an entourage V of X such that: \(\alpha) u_{0}^{-1}\) is uniformly continuous on \(\mathrm{V}(\mathrm{A})\) ; \(\beta)\) the relations \(u \in \Gamma\) and \((u(x_{0}), u(x)) \in \mathrm{V}\) for all \(x \in B\) together imply \(A \subset u(B)\) . Under these conditions, show that \(u \to u^{-1}\) (defined on \(\Gamma\) ) is continuous at \(u_{0}\) with respect to the topology of S-convergence on \(\Gamma\) .
  \end{enumerate}
\end{enumerate}

\begin{enumerate}
\def\labelenumi{\alph{enumi})}
\setcounter{enumi}{1}
\tightlist
\item
  Show that the condition \(\beta)\) may be replaced by the following: \(\beta'\) there exists a connected set \(\mathbf{C} \subset \mathbf{X}\) such that \(\mathbf{V}(\mathbf{A}) \subset \mathbf{C}\) and \(\mathbf{V}(\mathbf{C})\) is contained in the interior of \(u_0(\mathbf{B})\) . {[}Show that \(\beta') \Longrightarrow \beta)\) by reductio ad absurdum.{]}
\end{enumerate}

¶ 16) Let X be a uniform space and let \(\Gamma\) be the group of all automorphisms of the uniform structure of X.

\begin{enumerate}
\def\labelenumi{\alph{enumi})}
\item
  Show that the topology of uniform convergence on \(\Gamma\) is compatible with its group structure (use Exercises 12 and 15). Moreover, the uniformity induced on \(\Gamma\) by the uniformity of uniform convergence in X coincides with the right uniformity of the topological group \(\Gamma\) .
\item
  If X is the compact interval \([0, 1]\) of R, show that the topological group \(\Gamma\) defined in a) has no completion {[}construct a sequence of homeomorphisms \((u_{n})\) of X which is uniformly convergent but such that the sequence \((u_{n}^{-1})\) is not uniformly convergent{]}.
\item
  Show that if X is a complete Hausdorff uniform space, then the topological group \(\Gamma\) is complete with respect to its two-sided uniformity (Chapter III, § 3, Exercise 6). {[}Observe that if \(\Phi\) is a Cauchy filter on \(\Gamma\) with respect to this uniformity, then \(\Phi(x)\) and \(\Phi^{-1}(x)\) converge in X for all \(x \in X.\) {]}
\end{enumerate}

¶ 17) a) Show that if X is a locally compact, locally connected space, then the topology induced on the group \(\Gamma\) of all homeomorphisms of X by the compact-open topology coincides with the topology \(\mathcal{G}_{\beta}\) defined in no. 5 {[}use Exercise 15 b) and Proposition 12 of no. 5{]}.

\begin{enumerate}
\def\labelenumi{\alph{enumi})}
\setcounter{enumi}{1}
\tightlist
\item
  Let X be the locally compact subspace of R consisting of the points 0 and \(2^{n} (n \in \mathbf{Z})\) . If \(\Gamma\) is the group of all homeomorphisms of X, show that the topology induced on \(\Gamma\) by the compact-open topology is not compatible with the group structure of \(\Gamma\) .
\end{enumerate}

\begin{enumerate}
\def\labelenumi{\arabic{enumi})}
\setcounter{enumi}{17}
\tightlist
\item
  Let X be a locally compact space endowed with a uniformity compatible with its topology, and let \(\Gamma\) be the group of homeomorphisms of X. For each entourage V of X and each compact subset K of X, let \(G(\mathbf{K}, \mathbf{V})\) denote the set of pairs \((u, v)\) of homeomorphisms of X such that \((u(x), v(x)) \in \mathbf{V}\) and \((u^{-1}(x), v^{-1}(x)) \in \mathbf{V}\) for all \(x \in K\) .
\end{enumerate}

\begin{enumerate}
\def\labelenumi{\alph{enumi})}
\item
  Show that the sets \(G(\mathbf{K}, \mathbf{V})\) form a fundamental system of entourages of a uniformity U on \(\Gamma\) and that the topology induced by this uniformity is the topology \(T_{\beta}\) defined in no. 5.
\item
  Show that if X is complete with respect to its uniformity, then \(\Gamma\) is complete with respect to the uniformity U.
\item
  Show that \(\Gamma\) is complete with respect to the two-sided uniformity defined by the topology \(\mathfrak{G}_{\beta}\) {[}use Exercise 16 c){]}.
\item
  Take X to be the locally compact subspace of R consisting of points of the form \(n + 2^{-m}\) ( \(n \in Z, m\) an integer \(\geqslant 1\) ). Show that on the group \(\Gamma\) neither the uniformity U nor the uniformity of compact convergence is comparable with any of the three uniformities (left, right and two-sided) defined by the topology \(C_{\beta}\) .
\end{enumerate}

\begin{enumerate}
\def\labelenumi{\arabic{enumi})}
\setcounter{enumi}{18}
\tightlist
\item
  Let H be an equicontinuous group of homeomorphisms of a uniform space X; endowed with the topology of pointwise convergence, H is a topological group (no. 5, Corollary to Proposition 10).
\end{enumerate}

\begin{enumerate}
\def\labelenumi{\alph{enumi})}
\item
  Show that the left uniformity on H is finer than the uniformity of pointwise convergence and that these two uniformities coincide if H is uniformly equicontinuous.
\item
  Let u be the homeomorphism of the real line R defined by \(u(x) = x + 1\) for \(x \leqslant 0\) , \(u(x) = x + 1/(x + 1)\) for \(x \geqslant 0\) . Let H be the subgroup generated by u in the group of all homeomorphisms of R. Show that H is equicontinuous and is discrete with respect to the topology of pointwise convergence, but that the uniformity of pointwise convergence on H is not the discrete uniformity (observe that the limit \(u^{n+1}(x) - u^n(x)\) is 0 as n tends to \(+\infty\) ).
\item
  Take X to be the discrete space N of natural integers, endowed with the metric d such that \(d(m, n) = 1\) whenever \(m \neq n\) , and take H to be the group of isometries of N, which is uniformly equicontinuous. Show that the topological group obtained by endowing H with the topology of pointwise convergence has no completion {[}same method as in Exercise 16 b){]}.
\item
  Suppose that X is Hausdorff and complete and that H is uniformly equicontinuous. Show that the Cauchy filters on H with respect to the two-sided uniformity of the topological group \(H\) converge in the space \(\mathcal{F}_s(X;X)\) ; that the set \(H'\) of their limit points is a uniformly equicontinuous group of homeomorphisms of X, which (endowed with the topology of pointwise convergence) is complete with respect to its two-sided uniformity; and that H is dense in \(H'\) .
\end{enumerate}

\begin{enumerate}
\def\labelenumi{\arabic{enumi})}
\setcounter{enumi}{19}
\tightlist
\item
  Let X be the locally compact subspace of \(R^{2}\) consisting of the lines \(y = 0, y = 1/n\) ( \(n \geqslant 1\) ), and let G be the group of all homeomorphisms u of X whose restriction to each of the lines \(y = 0, y = 1/n\) is a translation of the form \((x, y) \to (x + a_{y}, y)\) . Show that G satisfies the conditions of Theorem 4 of no. 5, but that the topologies of compact and pointwise convergence on G are distinct.
\end{enumerate}

¶ 21) Let X be a locally compact space, and let T be a subset of C (X; X) consisting of surjective mappings. Let C be a topology on T which is finer than the topology of pointwise convergence and with respect to which T is locally compact, and consider the following property of the pair (T, C):

\begin{enumerate}
\def\labelenumi{(\Alph{enumi})}
\tightlist
\item
  For all \(u \in \mathbf{T}\) and \(v \in \mathbf{T}\) , we have \(u \circ v \in \mathbf{T}\) ; and for each \(u \in \mathbf{T}\) the mappings \(v \to u \circ v\) and \(v \to v \circ u\) of \(\mathbf{T}\) into itself are continuous with respect to \(\mathfrak{C}\) .
\end{enumerate}

\begin{enumerate}
\def\labelenumi{\alph{enumi})}
\item
  Suppose that X is compact and metrizable, that the pair (T, C) satisfies (A) and there exists a group G of homeomorphisms of X which is dense in T with respect to the topology C. Show that the set H of bijective mappings in T is the complement of a meagre set in T. {[}Note first, using Exercise 7, that every point of T has a compact metrizable neighbourhood (with respect to C). Let \((\mathbf{V}_{n})\) be a fundamental system of neighbourhoods in T of the identity element e of G, and let \(H_{n}\) be the set of all \(v \in T\) for which there exists \(u \in T\) such that \(v \circ u \in V_{n}\) ; observe that H contains the intersection of the sets \(H_{n}\) .{]}
\item
  Show that, under the hypotheses of \(a\) , the restriction to \(\mathbf{G} \times \mathbf{X}\) of the mapping \(\pi: (u, x) \to u(x)\) of \(\mathbf{T} \times \mathbf{X}\) into \(\mathbf{X}\) is continuous. {[}Show first that it is enough to prove that, for each \(x_0 \in \mathbf{X}\) , there exists \(u_0 \in \mathbf{H}\) such that \(\pi\) is continuous at \((u_0, x_0)\) , by establishing that this implies the continuity of \(\pi\) at \((e, x_0)\) . If \(\mathbf{V}\) is a compact metrizable neighbourhood of \(e\) in \(\mathbf{T}\) , show next, by using \(a\) ) and Chapter IX, § 5, Exercise 23, that there exists \(u_0 \in \mathbf{H} \cap \mathbf{V}\) such that \(\pi\) is continuous at \((u_0, x_0)\) .{]}
\end{enumerate}

\begin{enumerate}
\def\labelenumi{\arabic{enumi})}
\setcounter{enumi}{21}
\tightlist
\item
  Let X be a compact space, and let T be a subgroup of the group of all homeomorphisms of X, endowed with a locally compact topology \(\mathcal{C}\) for which condition (A) of Exercise 21 is satisfied. Let G be a countable subgroup of T, and let \(f\) be a continuous real-valued function on X. Consider the continuous mapping \(x \to \varphi(x)\) of X into \(\mathbf{I}^{\mathrm{G}}\) , where \(\mathbf{I} = f(\mathbf{X}) \subset \mathbf{R}\) and \(\varphi(x) = (f(u(x)))_{u \in \mathbf{G}}\) . Let Y be the image of X under \(\varphi\) ; Y is a compact metrizable subspace of \(\mathbf{I}^{\mathrm{G}}\) .
\end{enumerate}

\begin{enumerate}
\def\labelenumi{\alph{enumi})}
\item
  Show that the set of all \(v \in T\) such that \(\varphi \circ v = \varphi\) is a subgroup K of T whose normalizer in T contains G. Let R be the equivalence relation \(u \circ v^{-1} \in K\) on T, and let \(T'\) be the quotient space T/R and \(p: T \to T/R\) the canonical mapping. Show that \(T'\) is locally compact, and that if we set \(p(u)v = p(u \circ v)\) , then T operates on the right on \(T'\) in such a way that \(t' \to t'v\) is continuous on \(T'\) for all \(v \in T\) .
\item
  Let \(\mathbf{G}' = p(\mathbf{G})\) , let \(H'\) be the closure of \(G'\) in \(T'\) , and let \(\mathbf{H} = \overline{p}^{1}(\mathbf{H}')\) . Show that, if \(u \in H\) , the relation \(p(v) = p(v')\) implies \(p(u \circ v) = p(u \circ v')\) so that H operates on the left on \(T'\) ; moreover, if we set \(up(v) = p(u \circ v)\) for \(u \in H, v \in T\) , then the mapping \(t' \to ut'\) is continuous on \(T'\) . Finally, if \(u_1, u_2\) are elements of H such that \(p(u_1) = p(u_2)\) , we have \(u_1t' = u_2t'\) ; hence, by passing to the quotient, we define a mapping \((u', t') \to u't'\) of \(H' \times T'\) into \(T'\) , such that \(t' \to u't'\) is continuous on \(T'\) for all \(u' \in H'\) . Show that if \(u' \in H'\) and \(v' \in H'\) , then \(u'v' \in H'\) and that the law of composition so defined on \(H'\) induces a group structure on \(G'\) (with respect to which \(G'\) is isomorphic to GK/K).
\item
  If \(u \in \mathbf{H}\) , the relation \(\varphi(x) = \varphi(y)\) implies \(\varphi(u(x)) = \varphi(u(y))\) and consequently there is a unique mapping \(\tilde{u}: Y \to Y\) such that \(\varphi \circ u = \tilde{u} \circ \varphi\) . Show that \(\tilde{u}\) is continuous and surjective and that \(\tilde{u}_1 = \tilde{u}_2\) if and only if \(p(u_1) = p(u_2)\) ; this allows us to write \(\tilde{u} = \tilde{u}'\) , where \(u' = p(u)\) , and defines an injective mapping \(\psi: u' \to \tilde{u}'\) of \(H'\) into \(C_s(Y; Y)\) . Show that \(\psi\) is continuous and that \(\psi(u'v') = \psi(u') \circ \psi(v')\) ; using \(\psi\) to identify \(H'\) with a subset of \(C_s(Y; Y)\) and denoting by \(C'\) the topology on \(H'\) induced by that of \(T'\) , show that \((H', C')\) satisfies property (A) of Exercise 21.
\item
  Show that, if G is endowed with the topology induced by \(\mathfrak{G}\) , the mapping \((u, x) \to u(x)\) of \(\mathbf{G} \times \mathbf{X}\) into \(\mathbf{X}\) is continuous. {[}Show that, for every continuous real-valued function \(f\) on \(\mathbf{X}\) , the mapping \((u, x) \to f(u(x))\) is continuous on \(\mathbf{G} \times \mathbf{X}\) , by using \(c\) ) and Exercise 21; then apply Chapter IX, § 1, no. 5, Proposition 4.{]}
\end{enumerate}

\begin{enumerate}
\def\labelenumi{\arabic{enumi})}
\setcounter{enumi}{22}
\item
  Let X be a compact space and let T be a subset of \(\mathcal{C}(\mathbf{X};\mathbf{X})\) which is stable under the law of composition \((u,v)\to u\circ v\) . Suppose that T is endowed with a topology \(\mathcal{C}\) finer than that of pointwise convergence and with respect to which T is locally compact; also suppose that, for every countable stable subset S of T, the restriction to \(\mathbf{S}\times \mathbf{X}\) of the mapping \(\pi : (u,x)\to u(x)\) of \(\mathrm{T}\times \mathrm{X}\) into X is continuous. Show that, under these conditions, \(\pi\) is continuous on \(\mathrm{T}\times \mathrm{X}\) . {[}Let V be a compact subset of T (with respect to \(\mathcal{C}\) ); show that, for each countable subset D of V, the closure of D with respect to \(\mathcal{C}\) is contained in the closure of D with respect to the topology of uniform convergence, by using Corollary 1 of Theorem 3 (no. 4). Deduce that V is relatively compact in \(\mathcal{C}_u\) (X; X) (Chapter II, § 4, Exercise 6), and hence equicontinuous.{]}
\item
  Let X be a locally compact space, T a group of homeomorphisms of X, C a locally compact topology on T, which is finer than the topology of pointwise convergence and for which condition (A) of Exercise 21 is satisfied. Show that the mapping \((u, x) \to u(x)\) of \(T \times X\) into X is continuous with respect to C. {[}Extend the elements of T to homeomorphisms of the one-point compactification \(X'\) of X; then, by using Exercise 22 d), show that the result of Exercise 23 is applicable.{]}
\item
  Let G be a group endowed with a topology \(\mathcal{G}\) with respect to which G is locally compact and the translations \(t \to st\) and \(t \to ts\) are continuous on G for all \(s \in G\) . Show that \(\mathcal{G}\) is compatible with the group structure of G (Theorem of R. Ellis). {[}Identify G with the group of left translations of G, endowed with the topology of pointwise convergence, so that G is identified with a group of homeomorphisms of the one-point compactification X of G, endowed with the topology of pointwise convergence. Using Exercise 24, deduce that on G the topology of pointwise convergence in X is the same as the topology of uniform convergence in X (no. 4, Theorem 3, Corollary 1), and use Proposition 11 of no. 5.{]}
\end{enumerate}

¶ 26) a) Let X be a locally compact uniform space, G a group of homeomorphisms of X, and T the closure of G in \(\mathcal{C}_{s}\) (X; X). Show that T is compact with respect to the topology of pointwise convergence and is a group of homeomorphisms if and only if G is relatively compact in \(\mathcal{C}_{c}\) (X; X). {[}Use Exercises 12 and 24 and the Corollary to Theorem 4 (no. 5).{]}

\begin{enumerate}
\def\labelenumi{\alph{enumi})}
\setcounter{enumi}{1}
\tightlist
\item
  Let G be a locally compact, non-compact topological group and let X be its one-point compactification; G can be identified with a group of homeomorphisms of X. Show that G is not equicontinuous, although its closure in \(\mathcal{C}_{s}(\mathbf{X};\mathbf{X})\) is compact.
\end{enumerate}

\begin{enumerate}
\def\labelenumi{\arabic{enumi})}
\setcounter{enumi}{26}
\tightlist
\item
  Let X be a Hausdorff uniform space and let G be an equicontinuous group of homeomorphisms of X.
\end{enumerate}

\begin{enumerate}
\def\labelenumi{\alph{enumi})}
\item
  A point \(x_{0} \in X\) is said to be almost periodic with respect to G if the orbit of \(x_{0}\) with respect to G is relatively compact in X. Let Y be the (compact) closure of this orbit in X; then \(u(Y) \subset Y\) for all \(u \in G\) . Let \(\tilde{u}\) be the restriction of u to Y considered as an element of \(\mathcal{C}(Y; Y)\) ; show that \(\tilde{u}\) is a homeomorphism of Y onto itself and that the closure \(\Gamma\) in \(\mathcal{C}_{u}(Y; Y)\) of the image of G under the mapping \(u \to \tilde{u}\) is a compact group of homeomorphisms (no. 6, Theorem 4, Corollary) which is transitive on Y.
\item
  Suppose that X is complete. Then \(x_{0}\) is almost periodic with respect to G if and only if, for each neighbourhood V of \(x_{0}\) in X, there is a finite number of elements \(u_{i} \in G\) such that, for all \(u \in G\) , we have \(u_{i}^{-1}(u(x_{0})) \in V\) for at least one index i.
\item
  Suppose that \(x_{0}\) is almost periodic and that G is endowed with a topology compatible with its group structure. With the notation of a), the mapping \(u \to \tilde{u}\) of G into \(\Gamma\) (endowed with the topology of uniform convergence) is continuous if and only if, for each pair of distinct points x, y of Y, there is a neighbourhood U of e in G and a neighbourhood V of y such that the relation \(u \in U\) implies \(u(x) \notin V\) {[}use the fact that the topologies induced on \(\Gamma\) by those of \(\mathcal{C}_{u}(Y; Y)\) and \(\mathcal{C}_{s}(Y; Y)\) are the same{]}.
\end{enumerate}

\begin{enumerate}
\def\labelenumi{\arabic{enumi})}
\setcounter{enumi}{27}
\tightlist
\item
  Let G be a topological group and let X be the Banach space of bounded continuous mappings of G into C endowed with the norm
\end{enumerate}

\[
\left| \left| f \right| \right| = \sup _ {s \in \mathbf {G}} | f (s) |.
\]

For each \(f \in \mathbf{X}\) and each \(s \in \mathbf{G}\) , let \(U_s f\) denote the function \(t \to f(s^{-1}t)\) , which belongs to \(\mathbf{X}\) , so that the \(U_s\) , as \(s\) runs through \(\mathbf{G}\) , form a group of isometries \(\mathbf{G}'\) of \(\mathbf{X}\) . An element \(f \in \mathbf{X}\) is said to be an almost periodic function (on the left) on \(\mathbf{G}\) if \(f\) is an almost periodic element of \(\mathbf{X}\) with respect to the group \(\mathbf{G}'\) . For this to be so it is necessary and sufficient that, for each \(\varepsilon > 0\) , there should exist a finite number of elements \(s_i \in \mathbf{G}\) with the property that, for each \(s \in \mathbf{G}\) , there is at least one index \(i\) such that \(|f(s^{-1}t) - f(s_{i}^{-1}t)| \leqslant \varepsilon\) for all \(t \in \mathbf{G}\) .

\begin{enumerate}
\def\labelenumi{\alph{enumi})}
\tightlist
\item
  Suppose that \(f\) is almost periodic. Let \(\mathbf{Y}\) be the closure in \(\mathbf{X}\) of the orbit of \(f\) under \(\mathbf{G}'\) , and let \(\Gamma\) be the closure in \(\mathcal{C}_u(\mathbf{Y};\mathbf{Y})\) of the set of restrictions \(V_s\) of the \(U_s\) to \(\mathbf{Y} (s\in \mathbf{G})\) . For each \(\sigma \in \Gamma\) , let \(f_{\sigma}\in \mathbf{Y}\) be the image of \(f\) under \(\sigma\) . For each \(s\in \mathbf{G}\) , we have
\end{enumerate}

\[
f _ {\sigma} (t) = U _ {s} ^ {- 1} f _ {\sigma} (s ^ {- 1} t).
\]

Deduce that there is a continuous function \(\overline{f}\) on \(\Gamma\) such that \(f(s) = \overline{f}(V_s)\) {[}use Exercise 27 c){]}.

\begin{enumerate}
\def\labelenumi{\alph{enumi})}
\setcounter{enumi}{1}
\tightlist
\item
  Show by using a) that a function \(f \in X\) is almost periodic on G if and only if it is of the form \(g \circ \varphi\) , where \(\varphi\) is the canonical mapping of G into the compact group \(G^{c}\) associated with G (Set Theory, Chapter IV, § 3, no. 3, Example 8) and g is a continuous mapping of \(G^{c}\) into C.
\item
  Take G to be the additive group R. Show that if f is almost periodic on R, then for each \(\varepsilon > 0\) there exists a real number T \textgreater{} 0 such that every interval of R of length T contains a number s such that
\end{enumerate}

\[
\left| f (x + s) - f (x) \right| \leqslant \varepsilon \quad \text { for   all } \quad x \in \mathbf {R}
\]

(s is called an ``ε-period'' of f). {[}Use b) and Exercise 2 of Chapter V, § 1.{]}

\begin{enumerate}
\def\labelenumi{\arabic{enumi})}
\setcounter{enumi}{28}
\tightlist
\item
  Let X be a compact space and let G be a topological group operating continuously on X, such that every orbit of G is dense in X. Show that if, for every continuous real-valued function f on X, there exists \(x_{0} \in X\) such that \(s \to f(s.x_{0})\) is almost periodic on G, then G is equicontinuous (use the fact that X is homeomorphic to a closed subspace of a cube). Consider the converse, when G is abelian.
\end{enumerate}
% semantic-fragments: legacy-input-seam
\relax{}
